% Part X — The Geometry of Orchestration  (Slides 219 – 238)
\providecommand{\jugmark}[1]{\textbf{\color{jugeoBlue}#1}}

% ---------------------------------------------------------------------------
% Slide 219: Section divider
% ---------------------------------------------------------------------------
\begin{frame}
\frametitle{}
\vfill
\begin{center}
  {\huge\color{jugeoBlue}\textbf{Part X}}\\[16pt]
  {\LARGE The Geometry of Orchestration}\\[12pt]
  {\large\color{gray}\itshape How JuGeo coordinates proof, evidence, and repair}
\end{center}
\vfill
\end{frame}

% ---------------------------------------------------------------------------
% Slide 220: From Ideas to Action
% ---------------------------------------------------------------------------
\begin{frame}
\frametitle{From Ideas to Action}
\begin{center}
\begin{tikzpicture}[
  box/.style={draw, thick, rounded corners, minimum width=2.4cm,
              minimum height=0.9cm, font=\small, align=center},
  arr/.style={-{Stealth[length=3mm]}, thick}
]
  \node[box, fill=jugeoOrange!15] (idea) {Ideation\\discovers ideas};
  \node[box, fill=jugeoBlue!15, right=1.2cm of idea]
    (orch) {\jugmark{Orchestration}\\turns them into proofs};
  \node[box, fill=jugeoGreen!15, right=1.2cm of orch]
    (cert) {Certificate\\ships with code};

  \draw[arr, jugeoOrange] (idea) -- (orch);
  \draw[arr, jugeoGreen]  (orch) -- (cert);
\end{tikzpicture}
\end{center}
\vfill
\begin{itemize}
  \item Ideation generates \textbf{what} to prove
  \item Orchestration decides \textbf{how} to prove it:
    \begin{itemize}
      \item Which solver? Which evidence channel?
      \item How much budget? When to give up?
      \item What if two strategies disagree?
    \end{itemize}
  \item Think of it as an \jugmark{air traffic controller} for proof search
\end{itemize}
\end{frame}

% ---------------------------------------------------------------------------
% Slide 221: Analogy — Air Traffic Control
% ---------------------------------------------------------------------------
\begin{frame}
\frametitle{Analogy: Air Traffic Control}
\begin{columns}[T]
\column{0.48\textwidth}
\begin{block}{Air Traffic Control}
  \begin{itemize}
    \item Many planes (flights) at once
    \item Limited runway slots
    \item Must avoid collisions
    \item Priorities: emergencies first
    \item Reroute when weather changes
    \item Never lose track of anyone
  \end{itemize}
\end{block}
\column{0.48\textwidth}
\begin{exampleblock}{JuGeo Orchestration}
  \begin{itemize}
    \item Many proof tasks at once
    \item Limited compute budget
    \item Must avoid contradictions
    \item Priorities: high-trust first
    \item Reroute when solver stalls
    \item Never lose provenance
  \end{itemize}
\end{exampleblock}
\end{columns}
\vfill
\begin{center}
\jugmark{Orchestration is a control problem} ---
managing multiple concurrent activities under resource constraints,
with safety guarantees.
\end{center}
\end{frame}

% ---------------------------------------------------------------------------
% Slide 222: The Orchestration State
% ---------------------------------------------------------------------------
\begin{frame}
\frametitle{The Orchestration State}
\begin{center}
\begin{tikzpicture}[
  comp/.style={draw=jugeoBlue, fill=jugeoLight, rounded corners=3pt,
               minimum width=25mm, minimum height=8mm, font=\small\bfseries,
               line width=0.7pt, align=center},
]
  \node[comp] (cov) at (-3.5, 1.2) {Covers};
  \node[comp] (ctx) at (0, 1.2)    {Contexts};
  \node[comp] (sec) at (3.5, 1.2)  {Partial Sections};
  \node[comp] (tre) at (-3.5, 0)   {Treaties};
  \node[comp] (obl) at (0, 0)      {Obligations};
  \node[comp] (evi) at (3.5, 0)    {Evidence};
  \node[comp] (fro) at (-3.5,-1.2) {Frontier};
  \node[comp] (bud) at (0,-1.2)    {Budget};
  \node[comp] (pha) at (3.5,-1.2)  {Phase};

  \draw[jugeoBlue!30, thick] (-5.2, -2) rectangle (5.2, 2);
  \node[font=\scriptsize\bfseries\color{jugeoDark}] at (0, 2.3)
    {Orchestrator State};
\end{tikzpicture}
\end{center}
\vfill
At any moment, the orchestrator tracks:
\begin{itemize}
  \item What's \textbf{covered}, what's \textbf{partially proven},
        what \textbf{obligations} remain
  \item What \textbf{evidence} has been collected and at what \textbf{trust}
  \item How much \textbf{budget} is left, what \textbf{phase} we're in
  \item What \textbf{treaties} (agreements between subsystems) are active
\end{itemize}
\end{frame}

% ---------------------------------------------------------------------------
% Slide 223: Semantic Moves — The Available Actions
% ---------------------------------------------------------------------------
\begin{frame}
\frametitle{Semantic Moves --- The Available Actions}
\begin{center}
\small
\begin{tabular}{@{} l l @{}}
  \toprule
  \textbf{Move Category} & \textbf{Examples} \\
  \midrule
  \textcolor{jugeoBlue}{\textbf{Structural}}
    & Restrict to sub-coordinate, weaken a claim, split a goal \\
  \textcolor{jugeoGreen}{\textbf{Logical}}
    & Apply an inference rule, introduce a lemma, simplify \\
  \textcolor{jugeoOrange}{\textbf{Geometric}}
    & Refine a cover, attempt descent, compute overlaps \\
  \textcolor{jugeoRed}{\textbf{Evidence}}
    & Query Z3, run a test, ask the AI oracle \\
  \textcolor{jugeoDark}{\textbf{Treaty}}
    & Propose an interface agreement, challenge a claim \\
  \bottomrule
\end{tabular}
\end{center}
\vfill
\begin{itemize}
  \item At each step, the orchestrator has a \jugmark{menu of possible moves}
  \item Each move has a \textbf{precondition} (what must be true to apply it)
        and an \textbf{effect} (how it changes the state)
  \item The control law picks the \textbf{best move} given the current state,
        budget, and objectives
\end{itemize}
\end{frame}

% ---------------------------------------------------------------------------
% Slide 224: Fleet Competition — Let Strategies Compete
% ---------------------------------------------------------------------------
\begin{frame}
\frametitle{Fleet Competition --- Let Strategies Compete}
\begin{center}
\begin{tikzpicture}[
  member/.style={draw, thick, rounded corners, minimum width=2cm,
                 minimum height=0.8cm, font=\small, align=center},
  win/.style={member, fill=jugeoGreen!20, draw=jugeoGreen, line width=1.5pt},
  lose/.style={member, fill=gray!10, draw=gray},
  arr/.style={-{Stealth[length=3mm]}, thick}
]
  \node[lose] (a) at (0, 1.2)  {Strategy A\\``Use Z3''};
  \node[win]  (b) at (0, 0)    {Strategy B\\``Run tests''};
  \node[lose] (c) at (0, -1.2) {Strategy C\\``Ask AI''};

  \node[draw=jugeoGreen, thick, rounded corners, fill=jugeoGreen!10,
        minimum width=3cm, minimum height=0.9cm, font=\small\bfseries,
        align=center]
    (winner) at (5, 0) {Winner\\(best trust + cost)};

  \draw[arr, gray!50] (a.east) -- ++(1, 0) |- (winner.west);
  \draw[arr, jugeoGreen] (b.east) -- (winner.west);
  \draw[arr, gray!50] (c.east) -- ++(1, 0) |- (winner.west);
\end{tikzpicture}
\end{center}
\vfill
\begin{itemize}
  \item Multiple \jugmark{fleet members} propose solutions simultaneously
  \item Each bid includes: evidence, trust level, cost, and residual obligations
  \item A \textbf{bid evaluator} picks the winner using multi-criteria ranking
  \item Losers are archived as \textbf{fallback evidence} (not discarded)
  \item This is \textbf{competitive search}: the best approach wins on merit
\end{itemize}
\end{frame}

% ---------------------------------------------------------------------------
% Slide 225: Evidence Routing — Send It to the Right Place
% ---------------------------------------------------------------------------
\begin{frame}
\frametitle{Evidence Routing --- Send It to the Right Place}
\begin{center}
\begin{tikzpicture}[
  ch/.style={draw, thick, rounded corners, minimum width=2.2cm,
             minimum height=0.8cm, font=\small, align=center, fill=jugeoLight},
  router/.style={draw, thick, rounded corners=8pt, minimum width=2.6cm,
                 minimum height=1cm, font=\small\bfseries, align=center,
                 fill=jugeoBlue!15, draw=jugeoBlue, line width=1.2pt},
  arr/.style={-{Stealth[length=3mm]}, thick, jugeoBlue!50}
]
  \node[router] (r) at (0, 0) {Evidence\\Router};

  \node[ch] (z3) at (-4, 1.5)  {Z3 Solver};
  \node[ch] (rt) at (-4, 0)    {Runtime Tests};
  \node[ch] (ai) at (-4, -1.5) {AI Oracle};

  \node[ch] (obl) at (4, 0)    {Obligation\\to discharge};

  \draw[arr] (obl) -- (r);
  \draw[arr] (r) -- (z3);
  \draw[arr] (r) -- (rt);
  \draw[arr] (r) -- (ai);

  \node[font=\scriptsize\itshape, jugeoGreen] at (-2.2, 1.8) {arithmetic? $\to$ Z3};
  \node[font=\scriptsize\itshape, jugeoOrange] at (-2.2, -0.3) {IO behavior? $\to$ tests};
  \node[font=\scriptsize\itshape, jugeoRed] at (-2.2, -1.8) {heuristic? $\to$ AI};
\end{tikzpicture}
\end{center}
\vfill
\begin{itemize}
  \item Not every obligation should go to the same solver
  \item The \jugmark{evidence router} classifies each obligation and sends it
        to the best channel
  \item Routing is \textbf{trust-aware}: each channel has a trust ceiling
  \item If Z3 times out, the router \textbf{escalates} to the next channel
\end{itemize}
\end{frame}

% ---------------------------------------------------------------------------
% Slide 226: Trust-Aware Routing
% ---------------------------------------------------------------------------
\begin{frame}
\frametitle{Trust-Aware Routing}
\begin{center}
\begin{tikzpicture}[
  tier/.style={draw, thick, rounded corners, minimum width=7cm,
               minimum height=0.7cm, font=\small, align=center},
  arr/.style={-{Stealth[length=3mm]}, thick, jugeoRed!50}
]
  \node[tier, fill=jugeoGreen!15] (t1) at (0, 2.4)
    {\textsc{verified\_proof} --- Z3 with unsat core};
  \node[tier, fill=jugeoBlue!10] (t2) at (0, 1.2)
    {\textsc{solver\_discharged} --- Z3 without full proof};
  \node[tier, fill=jugeoOrange!10] (t3) at (0, 0)
    {\textsc{runtime\_witnessed} --- tests passed};
  \node[tier, fill=jugeoRed!8] (t4) at (0, -1.2)
    {\textsc{copilot\_suggested} --- AI proposed it};

  \draw[arr] (t1.south east) to[bend right=15]
    node[right, font=\scriptsize\itshape, text=jugeoRed]{escalation} (t3.north east);

  \node[font=\scriptsize, gray, align=center] at (5.2, 2.4) {ceiling: full proof};
  \node[font=\scriptsize, gray, align=center] at (5.2, 1.2) {ceiling: solver};
  \node[font=\scriptsize, gray, align=center] at (5.2, 0)   {ceiling: tests};
  \node[font=\scriptsize, gray, align=center] at (5.2, -1.2) {ceiling: suggestion};
\end{tikzpicture}
\end{center}
\vfill
\begin{itemize}
  \item Each evidence channel has a \jugmark{trust ceiling}
  \item AI cannot produce \textsc{verified\_proof}-level evidence, period
  \item When a channel fails, evidence is routed down the trust hierarchy
  \item The final trust level is \textbf{honest}: it reflects the weakest link
\end{itemize}
\end{frame}

% ---------------------------------------------------------------------------
% Slide 227: Treaty Negotiation — Finding Agreements
% ---------------------------------------------------------------------------
\begin{frame}
\frametitle{Treaty Negotiation --- Finding Agreements}
\begin{columns}[T]
\column{0.48\textwidth}
\begin{block}{The Problem}
Two subsystems make \textbf{overlapping claims} that don't quite agree.
\end{block}
\vspace{6pt}
\begin{exampleblock}{Example}
\begin{itemize}
  \item Module A: ``\texttt{process()} returns a list''
  \item Module B: ``\texttt{process()} returns a sorted list''
\end{itemize}
These aren't contradictions --- B is stronger than A.
Can they agree on an \textbf{interface}?
\end{exampleblock}
\column{0.48\textwidth}
\begin{center}
\begin{tikzpicture}[
  mod/.style={draw, thick, rounded corners, minimum width=2cm,
              minimum height=1cm, font=\small, align=center, fill=jugeoLight},
  treaty/.style={draw=jugeoGreen, thick, rounded corners, fill=jugeoGreen!15,
                 minimum width=2.8cm, minimum height=0.8cm, font=\small\bfseries,
                 align=center}
]
  \node[mod] (a) at (-1.5, 1.5) {Module A\\``returns list''};
  \node[mod] (b) at (1.5, 1.5)  {Module B\\``returns sorted list''};
  \node[treaty] (t) at (0, -0.2) {Treaty:\\``returns sorted list''};

  \draw[-{Stealth}, thick, jugeoBlue!50] (a.south) -- (t.north west);
  \draw[-{Stealth}, thick, jugeoBlue!50] (b.south) -- (t.north east);
\end{tikzpicture}
\end{center}
\end{column}
\end{columns}
\vfill
\begin{itemize}
  \item Treaties are \jugmark{negotiated interface agreements}
  \item They formalize how overlapping local claims combine
  \item Deadlocks are detected and broken by compromise strategies
\end{itemize}
\end{frame}

% ---------------------------------------------------------------------------
% Slide 228: Frontier Search — Exploring Proof Space
% ---------------------------------------------------------------------------
\begin{frame}
\frametitle{Frontier Search --- Exploring Proof Space}
\begin{center}
\begin{tikzpicture}[
  node/.style={draw, thick, circle, minimum size=8mm, inner sep=0pt,
               font=\tiny\bfseries},
  done/.style={node, fill=jugeoGreen!25},
  active/.style={node, fill=jugeoOrange!30, line width=1.5pt, draw=jugeoOrange},
  future/.style={node, fill=gray!10, draw=gray},
  edge/.style={thick, gray!30},
  hot/.style={thick, jugeoOrange!60}
]
  \node[done]   (a) at (0, 0)    {A};
  \node[done]   (b) at (1.8, 0.8){B};
  \node[done]   (c) at (1.8,-0.8){C};
  \node[active] (d) at (3.6, 1.4){D};
  \node[active] (e) at (3.6, 0)  {E};
  \node[active] (f) at (3.6,-1.4){F};
  \node[future] (g) at (5.4, 0.7){G};
  \node[future] (h) at (5.4,-0.7){H};

  \draw[edge] (a)--(b); \draw[edge] (a)--(c);
  \draw[hot]  (b)--(d); \draw[hot]  (b)--(e);
  \draw[hot]  (c)--(e); \draw[hot]  (c)--(f);
  \draw[edge] (d)--(g); \draw[edge] (e)--(g);
  \draw[edge] (e)--(h); \draw[edge] (f)--(h);
\end{tikzpicture}
\end{center}
\vfill
Each \jugmark{frontier node} carries predictions:
\begin{itemize}
  \item \textbf{Closure gain}: how much closer to a complete proof?
  \item \textbf{Stability}: will this move make things better or worse?
  \item \textbf{Cost}: time, solver queries, human attention
  \item \textbf{Diversity}: are we exploring enough different approaches?
\end{itemize}
The frontier scorer \textbf{ranks} nodes; the controller picks the best.
\end{frame}

% ---------------------------------------------------------------------------
% Slide 229: Phases — Exploration, Exploitation, Recovery
% ---------------------------------------------------------------------------
\begin{frame}
\frametitle{Phases --- Exploration, Exploitation, Recovery}
\begin{center}
\begin{tikzpicture}[
  phase/.style={draw, thick, rounded corners, minimum width=2.6cm,
                minimum height=1.1cm, font=\small, align=center},
  arr/.style={-{Stealth[length=3mm]}, thick, jugeoBlue!50}
]
  \node[phase, fill=jugeoBlue!15]
    (exp) at (0, 0) {\textbf{Exploration}\\Try many approaches\\Cast a wide net};
  \node[phase, fill=jugeoGreen!15]
    (opt) at (5, 0) {\textbf{Exploitation}\\Focus on the best\\Push to completion};
  \node[phase, fill=jugeoRed!10]
    (rec) at (2.5, -2.5) {\textbf{Recovery}\\Stalled? Backtrack\\Try something different};

  \draw[arr] (exp) -- node[above, font=\scriptsize]{found promising lead} (opt);
  \draw[arr] (opt) -- node[right, font=\scriptsize, text width=2cm, align=center]
    {stuck /\\budget low} (rec);
  \draw[arr] (rec) -- node[left, font=\scriptsize, text width=2cm, align=center]
    {new\\direction} (exp);
\end{tikzpicture}
\end{center}
\vfill
\begin{itemize}
  \item \jugmark{Phase transitions} are triggered by signals: progress rate,
        budget remaining, stall duration
  \item Like a chess engine switching from opening theory to endgame tactics
  \item The system detects \textbf{when} to switch, not just \textbf{what} to do
\end{itemize}
\end{frame}

% ---------------------------------------------------------------------------
% Slide 230: Backpressure — Knowing When to Pause
% ---------------------------------------------------------------------------
\begin{frame}
\frametitle{Backpressure --- Knowing When to Pause}
\begin{columns}[T]
\column{0.52\textwidth}
What happens when too many obstructions pile up?
\vspace{8pt}
\begin{alertblock}{Without backpressure}
  The system tries harder, creates more obligations, generates more
  obstructions --- a \textbf{death spiral} of wasted effort.
\end{alertblock}
\vspace{6pt}
\begin{exampleblock}{With backpressure}
  The system \textbf{pauses} proof production, \textbf{repairs} the
  worst obstructions, then \textbf{resumes} when healthy.
\end{exampleblock}
\column{0.44\textwidth}
\begin{center}
\begin{tikzpicture}[scale=0.9]
  \draw[-{Stealth}, thick] (0,0) -- (4,0) node[right, font=\scriptsize]{time};
  \draw[-{Stealth}, thick] (0,0) -- (0,3.5) node[above, font=\scriptsize]{obstructions};

  % Without backpressure
  \draw[jugeoRed, thick] plot[smooth]
    coordinates {(0,0.2) (0.5,0.5) (1,1) (1.5,1.8) (2,2.8) (2.5,3.3)};
  \node[font=\tiny, jugeoRed] at (2.8, 3.5) {no backpressure};

  % With backpressure
  \draw[jugeoGreen, thick] plot[smooth]
    coordinates {(0,0.2) (0.5,0.5) (1,1) (1.3,1.2) (1.5,0.8) (2,0.6) (2.5,0.4) (3,0.3) (3.5,0.2)};
  \node[font=\tiny, jugeoGreen] at (3.8, 0.4) {with backpressure};

  % Threshold
  \draw[jugeoOrange, dashed] (0, 1.3) -- (4, 1.3);
  \node[font=\tiny, jugeoOrange] at (4.3, 1.3) {threshold};
\end{tikzpicture}
\end{center}
\end{column}
\end{columns}
\end{frame}

% ---------------------------------------------------------------------------
% Slide 231: Budgets — Every Resource Is Finite
% ---------------------------------------------------------------------------
\begin{frame}
\frametitle{Budgets --- Every Resource Is Finite}
\begin{center}
\begin{tikzpicture}[
  dim/.style={draw, thick, rounded corners, minimum width=2.2cm,
              minimum height=0.7cm, font=\small, align=center, fill=jugeoLight}
]
  \node[dim] (t) at (0, 1.2)    {Time};
  \node[dim] (s) at (3, 1.2)    {Solver Queries};
  \node[dim] (c) at (6, 1.2)    {Copilot Tokens};
  \node[dim] (h) at (0, 0)      {Human Attention};
  \node[dim] (m) at (3, 0)      {Memory};
  \node[dim] (r) at (6, 0)      {Retries};
\end{tikzpicture}
\end{center}
\vfill
\begin{itemize}
  \item JuGeo tracks budgets across \jugmark{multiple dimensions}
  \item Invariant: \textbf{spent + reserved $\leq$ total} always holds
  \item Budget enforcement prevents runaway proof searches
  \item Allocation is \textbf{adaptive}: spend more on high-value tasks,
        less on diminishing returns
  \item Alerts fire when approaching limits; the controller adjusts strategy
\end{itemize}
\end{frame}

% ---------------------------------------------------------------------------
% Slide 232: Control Laws — How Decisions Are Made
% ---------------------------------------------------------------------------
\begin{frame}
\frametitle{Control Laws --- How Decisions Are Made}
\begin{center}
\small
\begin{tabular}{@{} l l l @{}}
  \toprule
  \textbf{Control Law} & \textbf{Strategy} & \textbf{When to use} \\
  \midrule
  \textbf{Greedy}      & Pick the locally best move   & Quick wins, simple goals \\
  \textbf{Lookahead}   & Simulate $k$ steps ahead     & Complex proofs, traps \\
  \textbf{Balanced}     & Spread effort across goals   & Many open obligations \\
  \textbf{Adaptive}    & Switch law based on feedback  & Large, varied tasks \\
  \bottomrule
\end{tabular}
\end{center}
\vfill
\begin{itemize}
  \item Each control law is a \jugmark{mathematical function} from
        state $\to$ move
  \item \textbf{Adaptive control} monitors progress and switches laws
        when the current one stalls
  \item This is formal control theory applied to proof search ---
        with convergence guarantees
\end{itemize}
\end{frame}

% ---------------------------------------------------------------------------
% Slide 233: Convergence — Are We Getting Closer?
% ---------------------------------------------------------------------------
\begin{frame}
\frametitle{Convergence --- Are We Getting Closer?}
\begin{columns}[T]
\column{0.52\textwidth}
JuGeo monitors \jugmark{convergence} in real time:
\begin{itemize}
  \item \textbf{Obligation count}: going down?
  \item \textbf{Evidence coverage}: going up?
  \item \textbf{Trust levels}: improving?
  \item \textbf{Obstruction count}: shrinking?
\end{itemize}
\vspace{6pt}
If all metrics improve $\to$ \textcolor{jugeoGreen}{\textbf{converging}}.\\
If metrics stall $\to$ \textcolor{jugeoOrange}{\textbf{phase switch}}.\\
If metrics worsen $\to$ \textcolor{jugeoRed}{\textbf{backtrack}}.
\column{0.45\textwidth}
\begin{center}
\begin{tikzpicture}[scale=0.85]
  \draw[-{Stealth}, thick] (0,0) -- (4.5,0) node[right, font=\scriptsize]{rounds};
  \draw[-{Stealth}, thick] (0,0) -- (0,3) node[above, font=\scriptsize]{obligations};

  \draw[jugeoGreen, thick] plot[smooth]
    coordinates {(0,2.5) (0.5,2.2) (1,1.8) (1.5,1.3) (2,0.9) (2.5,0.5) (3,0.3) (3.5,0.1)};

  \node[font=\tiny\bfseries, jugeoGreen] at (3.8, 0.4) {converging!};

  \draw[jugeoGreen, fill=jugeoGreen!30, opacity=0.3]
    plot[smooth] coordinates {(0,2.8) (0.5,2.5) (1,2.1) (1.5,1.6) (2,1.2) (2.5,0.8) (3,0.5) (3.5,0.3)}
    -- plot[smooth] coordinates {(3.5,0) (3,0.1) (2.5,0.3) (2,0.6) (1.5,1) (1,1.5) (0.5,1.9) (0,2.2)}
    -- cycle;
\end{tikzpicture}
\end{center}
\end{column}
\end{columns}
\end{frame}

% ---------------------------------------------------------------------------
% Slide 234: Synthesis Fallback — When Nothing Works
% ---------------------------------------------------------------------------
\begin{frame}
\frametitle{Synthesis Fallback --- When Nothing Works}
\begin{center}
\begin{tikzpicture}[
  box/.style={draw, thick, rounded corners, minimum width=2.6cm,
              minimum height=0.9cm, font=\small, align=center},
  arr/.style={-{Stealth[length=3mm]}, thick}
]
  \node[box, fill=jugeoRed!10] (stuck) {Proof search\\is stuck};
  \node[box, fill=jugeoOrange!15, right=1cm of stuck] (detect)
    {Deficit\\detected};
  \node[box, fill=jugeoBlue!15, right=1cm of detect]  (synth)
    {Synthesis\\campaign};
  \node[box, fill=jugeoGreen!15, right=1cm of synth]  (resume)
    {Resume with\\new theory};

  \draw[arr, jugeoRed!50]    (stuck)  -- (detect);
  \draw[arr, jugeoOrange!50] (detect) -- (synth);
  \draw[arr, jugeoGreen!50]  (synth)  -- (resume);
\end{tikzpicture}
\end{center}
\vfill
\begin{itemize}
  \item Sometimes the current theory is \textbf{insufficient} ---
        no combination of existing moves can close the proof
  \item The \jugmark{synthesis orchestrator} detects this ``theory deficit''
  \item It triggers a \textbf{synthesis campaign}: asks the ideation engine
        for new lemmas, analogies, or structural insights
  \item The new theory is injected and proof search resumes
  \item This is the \jugmark{ideation-orchestration loop}: each feeds the other
\end{itemize}
\end{frame}

% ---------------------------------------------------------------------------
% Slide 235: Treaty Memory — Learning from the Past
% ---------------------------------------------------------------------------
\begin{frame}
\frametitle{Treaty Memory --- Learning from the Past}
\begin{columns}[T]
\column{0.52\textwidth}
Every negotiation produces \jugmark{semantic capital}:
\begin{itemize}
  \item What interface agreements worked?
  \item What compromise strategies resolved deadlocks?
  \item What treaty patterns recur across projects?
\end{itemize}
\vspace{6pt}
JuGeo \textbf{archives} successful treaties so future verifications
can reuse them:
\begin{itemize}
  \item Faster convergence on similar codebases
  \item Accumulated institutional knowledge
  \item Replay and audit of past decisions
\end{itemize}
\column{0.44\textwidth}
\begin{center}
\begin{tikzpicture}[
  ep/.style={draw, thick, rounded corners, minimum width=2.4cm,
             minimum height=0.6cm, font=\scriptsize, align=center, fill=jugeoLight},
  arr/.style={-{Stealth[length=3pt]}, thick, jugeoBlue!40}
]
  \node[ep] (e1) at (0, 2.4)  {Episode 1: sorted list treaty};
  \node[ep] (e2) at (0, 1.6)  {Episode 2: auth scope treaty};
  \node[ep] (e3) at (0, 0.8)  {Episode 3: retry policy treaty};
  \node[draw=jugeoGreen, thick, rounded corners, fill=jugeoGreen!10,
        minimum width=2.8cm, minimum height=0.8cm, font=\scriptsize\bfseries,
        align=center]
    (lib) at (0, -0.4) {Treaty Library};

  \draw[arr] (e1) -- (lib);
  \draw[arr] (e2) -- (lib);
  \draw[arr] (e3) -- (lib);
\end{tikzpicture}
\end{center}
\end{column}
\end{columns}
\end{frame}

% ---------------------------------------------------------------------------
% Slide 236: The Full Loop — Ideation + Orchestration
% ---------------------------------------------------------------------------
\begin{frame}
\frametitle{The Full Loop}
\begin{center}
\begin{tikzpicture}[
  stage/.style={draw, thick, rounded corners, minimum width=2.4cm,
                minimum height=0.9cm, font=\small\bfseries, text=white,
                align=center},
  arr/.style={-{Stealth[length=3mm]}, thick, jugeoBlue!50}
]
  \node[stage, fill=jugeoOrange]                  (disc) at (0, 1.5)   {Discover};
  \node[stage, fill=jugeoBlue]                    (plan) at (3.5, 1.5) {Plan};
  \node[stage, fill=jugeoDark]                    (exec) at (7, 1.5)   {Execute};
  \node[stage, fill=jugeoGreen!80]                (cert) at (7, -0.5)  {Certify};
  \node[stage, fill=jugeoRed!70!black]            (rep)  at (3.5, -0.5){Repair};
  \node[stage, fill=jugeoOrange!80!black]         (feed) at (0, -0.5)  {Feed Back};

  \draw[arr] (disc) -- (plan);
  \draw[arr] (plan) -- (exec);
  \draw[arr] (exec) -- (cert);
  \draw[arr] (cert) -- (rep);
  \draw[arr] (rep)  -- (feed);
  \draw[arr] (feed) -- (disc);

  \node[font=\scriptsize\itshape, gray] at (1.75, 2) {ideation};
  \node[font=\scriptsize\itshape, gray] at (5.25, 2) {orchestration};
  \node[font=\scriptsize\itshape, gray] at (5.25, -1) {repair};
  \node[font=\scriptsize\itshape, gray] at (1.75, -1) {learning};
\end{tikzpicture}
\end{center}
\vfill
\begin{itemize}
  \item \textbf{Discover} $\to$ \textbf{Plan} $\to$ \textbf{Execute} $\to$
        \textbf{Certify} $\to$ \textbf{Repair} $\to$ \textbf{Feed Back}
        $\to$ repeat
  \item Ideation and orchestration form a \jugmark{continuous loop}
  \item Each cycle improves the system's knowledge and capability
  \item This is JuGeo's \textbf{cyclic maturity model}: the system gets
        smarter with every verification run
\end{itemize}
\end{frame}

% ---------------------------------------------------------------------------
% Slide 237: Orchestration vs. Traditional Proof Search
% ---------------------------------------------------------------------------
\begin{frame}
\frametitle{Orchestration vs.\ Traditional Proof Search}
\begin{center}
\small
\begin{tabular}{@{} l l l @{}}
  \toprule
  \textbf{Aspect} & \textbf{Traditional} & \textbf{JuGeo Orchestration} \\
  \midrule
  Strategy    & Fixed tactic script    & Adaptive control law \\
  Evidence    & Single channel (solver) & Multi-channel + routing \\
  Competition & None                   & Fleet bidding \\
  Failure     & Abort                  & Obstruction + repair \\
  Budget      & Timeout                & Multi-dimensional budget \\
  Memory      & None                   & Treaty archive \\
  Phases      & None                   & Explore / exploit / recover \\
  Feedback    & None                   & Synthesis fallback loop \\
  \bottomrule
\end{tabular}
\end{center}
\vfill
\begin{center}
Traditional proof search is a \textbf{single-threaded script}.\\
JuGeo orchestration is a \jugmark{multi-strategy, budget-aware, self-correcting control system}.
\end{center}
\end{frame}

% ---------------------------------------------------------------------------
% Slide 238: Orchestration Summary
% ---------------------------------------------------------------------------
\begin{frame}
\frametitle{Orchestration Summary}
\begin{center}
\small
\begin{tabular}{@{} r l @{}}
  \toprule
  \textbf{Component} & \textbf{What It Does} \\
  \midrule
  Semantic Moves    & Menu of available proof actions \\
  Fleet Competition & Multiple strategies compete on merit \\
  Evidence Routing  & Send obligations to the right channel \\
  Trust-Aware Routing & Respect trust ceilings per channel \\
  Treaty Negotiation & Agree on interfaces between subsystems \\
  Frontier Search   & Explore proof space systematically \\
  Phase Management  & Switch between exploration, exploitation, recovery \\
  Backpressure      & Pause when obstructions accumulate \\
  Budget Tracking   & Multi-dimensional resource management \\
  Control Laws      & Mathematical decision functions \\
  Convergence Monitor & Are we making progress? \\
  Synthesis Fallback & Create new theory when stuck \\
  Treaty Memory     & Learn from past negotiations \\
  \bottomrule
\end{tabular}
\end{center}
\vfill
\begin{center}
\large Next: \jugmark{Vision and Wrap-Up}
\end{center}
\end{frame}
